\documentclass[10pt]{article}
\usepackage[T1]{fontenc}
\usepackage[utf8]{inputenc}
\usepackage[french]{babel}
\usepackage{amsmath,amssymb,mathrsfs,mathpazo}
\usepackage{geometry}
\geometry{a4paper,margin=2cm}
\usepackage{graphicx,array,tabularx,tabularray,colortbl,multirow,multicol,booktabs,enumitem,pifont,ulem,fancybox,diagbox,cancel}
\usepackage{pstricks,pstricks-add,pst-plot,pst-tree,pst-node,pst-text,pst-func,pst-math,pst-bspline,pst-3dplot,pst-eucl,pst-solides3d}
\let\pscalculate\undefined
\usepackage{pst-bezier}
\usepackage{tikz,pgfplots,tkz-tab}
\usetikzlibrary{arrows,arrows.meta,patterns,calc,positioning,trees,decorations.markings,shapes.geometric,shapes.misc,shapes.arrows,intersections,decorations.pathmorphing,decorations.pathreplacing,backgrounds,shadows,babel}
\pgfplotsset{compat=1.18}
\usepgfplotslibrary{fillbetween}
\usepackage[np]{numprint}
\usepackage{listings,makecell,pas-tableur,twoopt}
\let\barre\undefined
\usepackage[pstricks]{tablvar}
\definecolor{darkgreen}{RGB}{0,100,0}
\providecommand{\red}{\color{red}}
\providecommand{\blue}{\color{blue}}
\providecommand{\green}{\color{green}}
\providecommand{\black}{\color{black}}
\providecommand{\white}{\color{white}}
\providecommand{\gray}{\color{gray}}
\providecommand{\cyan}{\color{cyan}}
\providecommand{\magenta}{\color{magenta}}
\providecommand{\yellow}{\color{yellow}}
\providecommand{\violet}{\color{violet}}
\providecommand{\orange}{\color{orange}}
\providecommand{\euro}{\texteuro}
\providecommand{\EUR}{\texteuro}
\providecommand{\ieme}{\textsuperscript{e}}
\providecommand{\degres}{\ensuremath{^\circ}}
\providecommand{\degree}{\ensuremath{^\circ}}
\providecommand{\vv}[1]{\overrightarrow{#1}}
\providecommand{\Coord}[1]{#1}
\providecommand{\decosix}{}
\providecommand{\starredbullet}{\ensuremath{\bullet}}
\pagestyle{empty}
\setlength{\parindent}{0pt}
\usepackage{adjustbox}\DeclareUnicodeCharacter{2236}{:}\DeclareUnicodeCharacter{2190}{\ensuremath{\leftarrow}}
\providecommand{\R}{\mathbb{R}}
\providecommand{\N}{\mathbb{N}}
\providecommand{\Z}{\mathbb{Z}}
\providecommand{\Q}{\mathbb{Q}}
\providecommand{\D}{\mathbb{D}}
\providecommand{\C}{\mathbb{C}}
\providecommand{\e}{{\mathrm{e}}}
\providecommand{\vect}[1]{\overrightarrow{#1}}
\providecommand{\vectt}[1]{\overrightarrow{\mathrm{#1}}}
\providecommand{\vvt}[1]{\overrightarrow{\mathrm{#1}}}
\providecommand{\np}[2][]{{#2~#1}}
\providecommand{\vv}[1]{{\overrightarrow{#1}}}
\renewcommand{\d}{{\mathrm{d}}}
\providecommand{\degres}{{^\circ}}
\providecommand{\degre}{{^\circ}}
\providecommand{\blue}{\color{blue}}
\providecommand{\red}{\color{red}}
\providecommand{\enskip}{\quad}
\providecommand{\llbracket}{\lbrack\!\lbrack}
\providecommand{\rrbracket}{\rbrack\!\rbrack}
\providecommand{\euro}{{\text{€}}}
\providecommand{\no}{{\text{n°}}}
\providecommand{\interval}[2]{{[#1;#2]}}
\providecommand{\intervFO}[2]{{[#1;#2[}}
\providecommand{\dotfill}{\ldots}

\providecommand{\ds}{\displaystyle}
\providecommand{\textding}[1]{#1}
\providecommand{\Oij}{\(\mathrm{O};\vec{\imath},\vec{\jmath}\)}
\providecommand{\Oijk}{\(\mathrm{O};\vec{\imath},\vec{\jmath},\vec{k}\)}
\providecommand{\Ouv}{\(\mathrm{O};\vec{u},\vec{v}\)}
\providecommand{\intervFF}[2]{\ensuremath{[#1;#2]}}
\renewcommand{\euro}{\ifmmode\text{\texteuro}\else\texteuro\fi}
\DeclareUnicodeCharacter{00B0}{\ifmmode^{\circ}\else\textdegree\fi}
\DeclareUnicodeCharacter{20AC}{\ifmmode\text{\texteuro}\else\texteuro\fi}
% Énoncés : Éducation nationale. Transcription et corrigés : APMEP, Denis Vergès et contributeurs. Découpage et classement : CoopMaths, Sébastien Lozano, Éric Elter et contributeurs. Adaptation de la présentation : Maths Explorer.
% CC BY-SA : https://coopmaths.fr/www/about/
% https://coopmaths.fr/alea/static/e3c/2020/tex/e3c_2020_00_sujet46_4.tex
\begin{document}
\noindent\textit{APMEP / CoopMaths. CC BY-SA. Présentation adaptée par Maths Explorer.}\par\medskip
\medskip

Des pucerons envahissent une roseraie.

On introduit alors des coccinelles, prédatrices des pucerons, à l'instant $t = 0$, et on s'intéresse à l'évolution du nombre de pucerons à partir de cet instant et sur une période de $20$ jours.

\bigskip

\textbf{Partie A :}

\medskip

Dans le repère ci-dessous, on a tracé :

\begin{itemize}[label=\textbullet]
\item La courbe $\mathcal{C}$ représentant le nombre de milliers de pucerons en fonction du nombre de jours écoulés depuis l'introduction des coccinelles.
\item La tangente $T$ à la courbe $\mathcal{C}$ au point d'abscisse 0 passe par les points A(0~;~2,1) et B(2~;~4,3).
\end{itemize}

\begin{center}
\psset{unit=0.6cm}
\begin{pspicture}(-1,-0.75)(22,7)
\psgrid[gridlabels=0pt,subgriddiv=5](0,0)(22,7)
\psaxes[linewidth=1.25pt,labelFontSize=\displaystyle]{->}(0,0)(0,0)(22,7)
\def\F{x 3 exp 0.003 mul x dup mul 0.12 mul sub 1.1 x mul add 2.1 add}
\psplot[plotpoints=2000,linewidth=1.25pt,linecolor=red]{0}{20}{x 3 exp 0.003 mul x dup mul 0.12 mul sub 1.1 x mul add 2.1 add}
\uput[ur](16,1.25){\red $\mathcal{C}$}\uput[u](20,0){\small $t$, nombre de jours}\uput[r](0,6.7){\small  $f(t)$ nombre de pucerons en milliers}

\psline[linewidth=1.25pt]{->}(0,2.1)(2,4.3)
\uput[ul](2,4.3){$T$}
\end{pspicture}
\end{center}

\medskip

\begin{enumerate}
\item Déterminer par lecture graphique le nombre de pucerons à l'instant où l'on introduit les coccinelles puis le nombre maximal de pucerons sur la période de $20$ jours.
\item On assimile la vitesse de prolifération des pucerons à l'instant $t$ au nombre dérivé $f'(t)$.

Déterminer graphiquement la vitesse de prolifération des pucerons à l'instant $t = 0$.
\end{enumerate}

\textbf{Partie B :}

\medskip

On modélise l'évolution du nombre de pucerons par la fonction $f$ définie, pour tout $t$ appartenant à l'intervalle [0~;~20], par:

\[f(t) = 0,003t^3 - 0,12t^2 + 1,1t + 2,1\]

où $t$ représente le nombre de jours écoulés depuis l'introduction des coccinelles et
$f(t)$ le nombre de pucerons en milliers.

\medskip

\begin{enumerate}
\item Déterminer $f'(t)$ pour tout $t$ appartenant à l'intervalle [0~;~20] où $f'$ désigne la
dérivée de la fonction $f$.
\item Dresser le tableau de signes de $f'(t)$ sur l'intervalle [0~;~20].
\item En déduire le tableau des variations de la fonction $f$ sur l'intervalle [0~;~20]. Préciser les images des valeurs de $t$ apparaissant dans le tableau.
\end{enumerate}

\end{document}
